% TOP_PROBLEM: {"id": 4600016, "title": "Commuting powers conjecture", "category_id": 11, "category": {"id": 11, "name": "dynamical_systems", "display_name": "Dynamical Systems", "description": "Problems about long-term behavior of deterministic systems, Hamiltonian dynamics, and periodic orbits.", "slug": "dynamical-systems", "order_index": 11, "created_at": "2026-07-31T15:26:25.671Z"}, "status": "open", "problem_number": "AMR-045-0016", "set_id": 13, "statusQualification": "The standard nontrivial mixing-SFT convention is explicit; including a one-point system paired with an infinite system would make the statement false."}
% FIRST_POSTED: 2026-10-01
% ENTRY_KIND: statement_only
% UnsolvedMath ID 4600016; source catalogue code AMR-045-0016
% !TeX program = lualatex
% Definition and sources only; this file is not an LLM solution attempt.
\documentclass[11pt,a4paper]{article}
\usepackage[margin=25mm]{geometry}
\usepackage{fontspec}
\setmainfont{lmroman10-regular.otf}[BoldFont=lmroman10-bold.otf,
 ItalicFont=lmroman10-italic.otf,BoldItalicFont=lmroman10-bolditalic.otf]
\usepackage{amsmath,amssymb,mathtools}
\usepackage{unicode-math}
\setmathfont{latinmodern-math.otf}
\AtBeginDocument{\let\setminus\smallsetminus}
\usepackage{xurl,hyperref}
\hypersetup{colorlinks=true,urlcolor=blue}
\setcounter{secnumdepth}{0}
\setlength{\parindent}{0pt}
\setlength{\parskip}{5pt}
\setlength{\emergencystretch}{3em}
\begin{document}
\section*{Commuting powers conjecture}\label{4600016}
{\small UnsolvedMath ID 4600016 · AMR-045-0016 · Dynamical Systems · AMR Open Problem Lists}
\subsection{Definitions and mathematical statement}
Let $(X,S)$ and $(Y,T)$ be nontrivial mixing two-sided shifts of finite
type over finite alphabets. Nontrivial here excludes the one-point system;
in this setting the spaces are Cantor spaces and have positive entropy.
Mixing means that for any nonempty open sets $U,V$, every sufficiently
large iterate of $U$ meets $V$.

Say that two homeomorphisms \emph{can commute} if they have topologically
conjugate representatives on a common compact space which commute.
The commuting powers conjecture asserts that there is $N=N(S,T)$ such
that for every pair of integers $i,j\ge N$ there exists a homeomorphism
$H_{i,j}:Y\to X$ for which
\[
 S^i\bigl(H_{i,j}T^jH_{i,j}^{-1}\bigr)
   =\bigl(H_{i,j}T^jH_{i,j}^{-1}\bigr)S^i.
\]
Equivalently, $S^i$ and $T^j$ can be realized as two commuting generators
of a $\mathbb Z^2$ action, each retaining its specified conjugacy class.

The same threshold must work for every sufficiently large pair, but the
identification $H_{i,j}$ may depend on that pair. No common identification
for all powers is requested. The nontrivial convention is necessary:
a one-point mixing shift cannot be placed on the same underlying space
as an infinite one. The conjecture concerns two-sided shifts, with no
requirement that their defining adjacency matrices commute or have equal
entropy.
\subsection{Short English statement}
Can every pair of nontrivial mixing finite-type shifts be made to commute after taking any sufficiently large powers?
\subsection{Sources}
\begin{itemize}
\item[S1] ULAM.AI and the UnsolvedMath Contributors, supplied problems.json, record 4600016 (AMR-045-0016), AMR Open Problem Lists. Catalogue identity and source transcription; CC BY 4.0.
\url{https://huggingface.co/datasets/ulamai/UnsolvedMath}.
\item[S2] Boyle - Open Problems in Symbolic Dynamics (2008), Problem/Question/Conjecture 13.1. Original source indexed by AMR-045-0016.
\url{https://www.math.umd.edu/~mboyle/papers/openfinalsub3nov2007.pdf}.
\end{itemize}
\end{document}
